\section{Implementation}
\label{sec:implementation}

% Data generation, scaling/centering decision (with your critical justification), train/test split
% Code architecture (high-level: classes/functions, the repo)
% Pseudocode for non-trivial algorithms (GD variants, bootstrap, CV)
% Libraries (numpy, scikit-learn, JAX/autograd)

The problem solution was implemented as a Python project. This section aims to introduce this implementation, as well as decisions/trade-offs made and show how it was verified. The source code, exploratory notebooks, tests, and experimental scripts required to reproduce the results are available in the publicly hosted project repository~\cite{risethFYSSTK4155Project12026}.

\subsection{Data generation}

We generate noisy observations from Runge's function by adding i.i.d. Gaussian noise, as shown in \cref{eq:data_generating_process}, where $f$ is Runge's function (\ref{eq:runges_function}) and the noise $\bm{\epsilon}$ has standard deviation $\sigma=0.1$. We generate the data points $(x_i, y_i)$ by drawing $n$ $x$-values uniformly from the interval $[-1, 1]$, evaluating $f$ at each, and adding the noise. An example of the resulting data points is shown in \cref{fig:runge_data}.

\begin{figure}[htbp]
    \centering
    \includegraphics[width=\linewidth]{figures/data/runge_data.pdf}
    \caption{Noisy samples ($n=100$, $\sigma=0.1$, seed 42) drawn from Runge's function, shown against the noise-free ground truth.}
    \label{fig:runge_data}
\end{figure}

The data was subsequently split into training and test sets (80/20), before standardization. The split was performed before standardization to avoid data leakage. Standardization was applied using \texttt{StandardScaler}, so that each column of the design matrix has mean zero and unit standard deviation, and centered target has mean zero. This improves the conditioning of the optimization landscape and ensures that the penalties in Ridge and Lasso regression treat each coefficient fairly, since the features are on the same scale~\cite{hjorth-jensenMachineLearningPhysical2026, geronHandsonMachineLearning2023}.

% "The StandardScaler in the pipeline is not decoration. A degree-six polynomial basis built from x ∼ N (0, 1) has columns whose scales differ by orders of magnitude, so the design matrix is severely ill-conditioned in the sense of Eq. (1.77); without scaling, the cross-validation curve is dominated by numerical noise at the tails and its minimum lands at whichever end of the grid one happened to choose. This is the Vandermonde conditioning problem of Chapter 1 appearing in statistical clothing, and it is a good illustration of why the two halves of this book cannot be kept apart." -- hjorth-jensenMachineLearningPhysical2026

\subsection{Code architecture}

The source code was implemented in Python using the \texttt{numpy}, \texttt{sklearn}, and \texttt{jax} libraries~\cite{harrisArrayProgrammingNumPy2020, pedregosaScikitlearnMachineLearning2011,bradburyJAXComposableTransformations2018}. For \texttt{jax}, the precision was changed to double precision. % Double precision was a requirement from the course lecturer Hjorth-Jensen; is it worth mentioning? Nevertheless, citing the libraries used was a requirement.

\subsubsection{Regression}

The regression models all inherit from the \texttt{LinearModel} abstract base class, which follows the \texttt{sklearn} estimator convention with \texttt{fit} and \texttt{predict} functions, and hyperparameters set in the constructor. This lets our estimators be dropped directly into \texttt{sklearn} utilities such as \texttt{KFold} and \texttt{cross\_validate}. The base class provides the shared \texttt{predict} method and input validation, while each subclass implements \texttt{fit}. OLS solves the least-squares problem via the SVD-based pseudoinverse, which remains stable for ill-conditioned or rank-deficient design matrices. Ridge instead solves the normal equations from \cref{eq:ridge-theta-hat} directly with \texttt{numpy.linalg.solve}, since the penalty term guarantees a well-conditioned system, with the intercept optionally excluded from the penalty. Lasso, which lacks a closed-form solution, is fit via gradient descent. % I'm unsure how relevant this really is to introduce.

\subsubsection{Resampling}

Divided into \texttt{bootstrap.py} and \texttt{cross\_validation.py}, where bootstrap uses numpy and the cross-validation uses sklearn's \texttt{KFold} class and \texttt{cross\_val\_score} function.

\subsubsection{Gradient descent}

The gradient descent code is implemented as its own package (\texttt{optimization}) containing all the different variants of gradient descent. They all roughly follow \cref{alg:gd_basic}, but replace line 3 with their own corresponding update equation.

\begin{algorithm}
    \caption{Basic gradient descent}
    \label{alg:gd_basic}
    {\small
    \begin{algorithmic}[1]
        \Require Initial parameters $\bm{\theta}^{(0)}$, learning rate $\eta > 0$, tolerance $\varepsilon > 0$, maximum iterations $N_{\max}$
        \Ensure Approximate minimizer $\bm{\theta}^*$ of $J(\bm{\theta})$
        \State $k \gets 0$
        \While{$\|\nabla J(\bm{\theta}^{(k)})\| > \varepsilon$ \textbf{and} $k < N_{\max}$}
            \State $\bm{\theta}^{(k+1)}
                \gets \bm{\theta}^{(k)}
                - \eta \nabla J(\bm{\theta}^{(k)})$ \Comment{\Cref{eq:gd_basic}}
            \State $k \gets k + 1$
        \EndWhile
        \State \Return $\bm{\theta}^{(k)}$
    \end{algorithmic}
    }
\end{algorithm}

\subsection{Verification}

% Reproduce known results: GD converges to closed-form OLS/Ridge; analytical vs autodiff gradients agree to machine precision; your k-fold matches scikit-learn on identical folds; learning-rate bound vs Hessian eigenvalue. This is where you show the code is correct before trusting its output.

% Should we briefly mention that all code that has been produced (in src/) have their own associated tests? For example does Ridge have a test that ensures lambda=0 matches the result of OLS.
